\subsection{The separate two-object frontier at nineteen arrows}

\begin{proposition}[Proved: strong two-object support through nineteen]
\label{CP2:A-strong-two-nineteen}
A strongly connected two-object category with at most nineteen
arrows cannot have classifying space of type $\mathbb CP^2$.
\end{proposition}
\begin{proof}
An object retract can be deleted, leaving a monoid with fewer than
twenty-three elements; the complete monoid minimum excludes it.
Otherwise both identities lie outside the categorical minimal ideal.
Use the notation of Lemma~\ref{CP2:A-high-Rees}: total row and column
counts $r,c$, and minimal local group $G$.  If either total is two,
all local counts on that side equal one, so the classifying space
is that of a finite group and has no positive rational homology.
Hence $r,c\ge3$ and
\[
 |G|rc+2\le19.
\]
This forces $G=1$.  If one total is three, its two positive local
counts are one and two.  The binary-representable obstruction of
Lemma~\ref{CP2:A-baseline-binary-representable} kills the rational
second products.  Otherwise $r,c\ge4$, and the only possible pair
is $r=c=4$.

Put $R=\mathbb QC$ and $A=R/\mathbb QK$.  The sixteen arrows in the
minimal band leave $\dim A\le3$.  Both object identities survive in
$A$ as nonzero orthogonal idempotents whose sum is the identity.
Thus $A$ is nonlocal.  The categorical band--Tor formula gives
\[
 H_j(BC;\mathbb Q)=\operatorname{Tor}_{j-2}^{A}(W,V),\qquad j\ge2.
\]
A projective plane would have nonzero second Tor and bounded higher
Tor.  Lemma~\ref{A:nonlocal-four} excludes this for every nonlocal
algebra of dimension at most four, in particular for this $A$.
The argument includes the pure eighteen-arrow band case as well.
\end{proof}

\begin{lemma}[Proved: all small sources opposite the right-zero corner]
\label{CP2:A-eight-source-right-zero-all-states}
For $|M|\le8$, every bounded-homology collage over a finite
left-$R_3$, right-$M$ biset has zero rational degree-two products.
There is no state-count restriction.
\end{lemma}
\begin{proof}
If $H^2(BM;\mathbb Q)=0$, both corner restrictions vanish and the
relative suspension quotient kills the products.  Otherwise the
minimal ideal has at least two rows and two columns, by rational
minimal-group collapse and the balanced tensor formula.  Its group
must be trivial: a nontrivial group already requires at least
$2\cdot2\cdot2+1=9$ elements including the outside identity.
The only possible minimal-band shapes are $2\times2$, $2\times3$
and $3\times2$.  The first two use the general all-state actor
quotient theorems.  In the last case there are at most two outside
elements, so Theorem~\ref{CP2:A-two-column-two-outside-graph} applies
to the original bounded collage.  No boundedness is transferred
through an actor quotient in this argument.
\end{proof}

\begin{lemma}[Proved: the twelve-by-three four-state endpoint]
\label{CP2:A-twelve-three-four-nineteen}
There is no nineteen-arrow projective-plane collage of profile
$(12,3,4)$.
\end{lemma}
\begin{proof}
For $N=R_3$, Proposition~\ref{CP2:A-right-zero-four-states} gives an
actual ordinary act-cone projective plane with the same source and
at most four states.  It has at most seventeen arrows, impossible.
Every nongroup three-element monoid other than $R_3$ has a
nonidentity minimal idempotent $e$ with $eN=eNe$: its proper minimal
ideal has one column.  The sink version of
Theorem~\ref{CP2:A-extreme-idempotent-corner-reduction} replaces that
corner by $eNe$ and the states by $eS$, by an actual homotopy
equivalence.  At least one arrow is removed, contradicting the
complete eighteen-arrow lower bound.

The remaining monoid is $C_3$.  If its action is nontrivial, its
four-state orbit set has at most two elements.  Exact rational
group coinvariants identify the collage's homology and ring with
that ordinary act cone, whose second products vanish by the
unconditional one- and two-state obstructions.  If the action is
trivial, the ordinary act cone is an actual categorical retract.
The collapse is a rational homology equivalence because $BC_3$ is
rationally acyclic.  The split integral homology, trivial fundamental
group and primitive square therefore make the retract an actual
projective plane.  Its smaller arrow budget is again impossible.
\end{proof}

\begin{proposition}[Proved: an initial nineteen-arrow collage inventory]
\label{CP2:A-two-object-nineteen-initial-scope}
After the chosen positive-source, fixed-point-free orientation,
the initial nineteen-arrow frontier consists of thirty-two ordered
profiles.  Four have a corner of order at most two:
\[
 (12,1,6),\quad(13,1,5),\quad(12,2,5),\quad(13,2,4).
\]
The other twenty-eight have $m\ge5$, $n\ge3$, $s\ge4$ and
$m+n+s=19$.  They are exactly all thirty-six triples satisfying
those inequalities except
\[
\begin{gathered}
 (5,3,11),\ (6,3,10),\ (7,3,9),\ (8,3,8),\ (12,3,4),\\
 (5,9,5),\ (5,10,4),\ (6,9,4).
\end{gathered}
\]
This is a scope reduction only; no nineteen-arrow finite search is
claimed complete here.
\end{proposition}
\begin{proof}
The complete eighteen-arrow theorem and the monoid minimum reduce
the question to an exactly nineteen-arrow connected two-object
category.  Strong support is excluded above, so the category is a
collage.  The positive-source orientation gives $m\ge5$ and a
right act with no global fixed state.  There are
$\binom{14}{2}=91$ positive triples $m+n+s=19$ before further
reductions.  At most three connecting arrows are excluded by the
universal one-, two- and three-cross results.

For $n\le2$ and $m\le11$, the bounded-eleven act theorem, with
exact orbit comparison when needed, applies.  For $n=1,m=14,s=4$
use the four-state theorem through order fourteen.  The four
displayed small-corner profiles are the remaining ones; this
statement leaves their rational-versus-actual comparison questions
open.  In their order-two corners a nontrivial idempotent gives an
actual proper corner reduction, so only the group $C_2$ with a
nontrivial action needs further work.

For $n\ge3,s\ge4$ the count is $\binom9{2}=36$.
The four profiles with $n=3,m\le8$ use the preceding all-state
right-zero lemma; every other three-element corner has projective
trivial right rational module, so bounded-eleven exact orbit
comparison removes it.  The profile $(12,3,4)$ is the preceding
endpoint lemma.  The two $m=5,s\le5$ profiles use the arbitrary
opposite-corner $P_5$ theorem in the chosen orientation.  Finally
$(6,9,4)$ uses the oriented six-source four-state theorem, whose
opposite corner is unrestricted.  These eight profiles are
distinct, leaving twenty-eight, and therefore thirty-two in all.
The fresh executable inventory records every one of the ninety-one
triples and each applied proof label.
\end{proof}

\paragraph{Computational priority and exact limits.}
The first finite family selected for investigation is the pure
$2\times4$ band with an identity, opposite $R_3$, on seven states:
it belongs to $(9,3,7)$ and is immediately beyond the proved
six-state rank-one-top obstruction.  The general $2\times2$ and
$2\times3$ opposite-$R_3$ proofs are used before considering a new
order-eleven source census.  A newly induced quotient is used for
products only under its stated theorem; its higher homology is not
assumed bounded merely because it preserves degree two.
